\section{Introdução}
\label{sec:educacao-intro}

A educação odontológica contemporânea encontra-se no epicentro de uma transição paradigmática, impulsionada por um conjunto complexo de desafios que tensionam, de maneira irrevogável, o modelo tradicional de ensino. Historicamente embasada no preceptorado clínico\footnote{O preceptorado clínico refere-se à tutoria direta exercida por um profissional experiente (preceptor) sobre um estudante em ambiente clínico real. Embora essencial, apresenta limitações estruturais de escala, dependência de casos incidentais e variabilidade na padronização das avaliações formativas e somativas.} e na prática laboratorial em manequins estáticos ou dentes extraídos, a pedagogia odontológica enfrenta a necessidade premente de um treinamento prático supervisionado de alta densidade, o qual demanda laboratórios pré-clínicos onerosos, insumos de uso único e milhares de horas de tutoria individualizada. Concomitantemente, as instituições de ensino superior lidam com restrições orçamentárias crescentes, dilemas éticos relacionados ao uso de pacientes reais nas fases iniciais de aprendizado e a pressão indomável por matrizes curriculares que incorporem competências digitais avançadas desde a base estrutural do curso~\cite{duggal2026exploring}.

\destaque{A transição da odontologia analógica para a Odontologia de Precisão não exige apenas novos equipamentos, mas uma reformulação epistemológica de como o conhecimento clínico e motor é adquirido e internalizado pelo aprendiz.}

Neste contexto de inflexão, os gêmeos digitais (\textit{Digital Twins}) emergem não apenas como ferramentas acessórias, mas como vetores de uma reestruturação fundamental do espaço educacional, capazes de conciliar demandas aparentemente ortogonais. Ao oferecerem representações virtuais dinâmicas, matematicamente rigorosas e fisiologicamente responsivas de pacientes, tecidos, procedimentos e instrumentais, eles permitem que o estudante transite do estágio cognitivo ao autônomo\footnote{O modelo de Fitts e Posner (1967) divide o aprendizado motor em três fases sucessivas: 1) Estágio cognitivo (compreensão intelectual do movimento e alta dependência de atenção consciente); 2) Estágio associativo (refinamento das respostas motoras e detecção de erros próprios); 3) Estágio autônomo (execução automática com mínima carga cognitiva, permitindo foco em decisões clínicas estratégicas).} na aquisição de habilidades motoras em um ambiente imersivo, estritamente seguro e parametrizado~\cite{katsoulakis2024digital}. Diferentemente das simulações estáticas, que se limitam à morfologia macroscópica, os gêmeos digitais na educação odontológica incorporam modelagem de tecidos moles e duros baseada em física, variabilidade anatômica estocástica, resposta biomecânica a forças de corte e cisalhamento, e evolução temporal das patologias, aproximando a carga cognitiva da experiência educacional à da realidade clínica irrestrita~\cite{roy2025applications}.

A virtualização do paciente desvincula o treinamento da dependência de casos clínicos incidentais nas clínicas-escola, garantindo que todos os discentes sejam expostos a um espectro padronizado de complexidades, desde anatomias canônicas até variações morfológicas raras e anomalias de desenvolvimento~\cite{khoshfekr2025review}. Este capítulo destrincha o estado da arte das tecnologias educacionais baseadas em gêmeos digitais, examinando detalhadamente a arquitetura dos simuladores com \textit{feedback} háptico\footnote{O feedback háptico (do grego \textit{haptikós}, relativo ao tato) é a tecnologia que replica as sensações de toque, resistência mecânica, textura e vibração ao usuário através de forças contrárias geradas por motores eletromecânicos. No simulador odontológico, calcula em tempo real o produto da rigidez da superfície virtual pela penetração da broca virtual, gerando torque físico real na mão do estudante.}, as evidências epistemológicas e cognitivas do aprendizado baseado em simulação (SBL), as topologias de plataformas distribuídas que viabilizam a tele-educação e, de forma imperativa, a contribuição específica do ecossistema OpenCode para a consolidação de material didático escalável, revisado por pares e ancorado em ontologias padronizadas. Ao final, aprofundam-se as barreiras sociotécnicas e institucionais que ainda mitigam a adoção ubiquitária dessas tecnologias no hemisfério acadêmico.
